\subsection{BRACKET IN THE LIE ALGEBRA}

For \(\alpha \in \mathbf{N}^n\) , let \(e_{\alpha}\) be the point distribution \(\frac{1}{\alpha!} \frac{\partial^{\alpha}}{\partial x^{\alpha}}\) at the origin. In particular, \(e_k = e_{\varepsilon_k} = \frac{\partial}{\partial x_k}\) . The \(e_{\alpha}\) form a basis of the vector space \(\mathrm{U}(G)\) . If \(f\) is an analytic function on an open neighbourhood of 0 in G and \(f(x) = \sum_{\alpha} \lambda_{\alpha} x^{\alpha}\) is its expansion as an integral series about the origin, then

In particular,

\[
\langle e _ {\alpha}, f \rangle = \lambda_ {\alpha}.
\]

\[
\langle e _ {\alpha}, x ^ {\gamma} \rangle = \delta_ {\alpha \gamma}.
\]

Hence

\[
\begin{array}{r l} \langle e _ {\alpha} * e _ {\beta}, x ^ {\gamma} \rangle & = \langle e _ {\alpha} \otimes e _ {\beta}, (x. y) ^ {\gamma} \rangle \\ & = \langle e _ {\alpha} \otimes e _ {\beta}, \sum_ {\alpha^ {\prime}, \beta^ {\prime}} c _ {\alpha^ {\prime} \beta^ {\prime} \gamma} x ^ {\alpha^ {\prime}} y ^ {\beta^ {\prime}} \rangle \\ & = \sum_ {\alpha^ {\prime}, \beta^ {\prime}} c _ {\alpha^ {\prime} \beta^ {\prime} \gamma} \langle e _ {\alpha}, x ^ {\alpha^ {\prime}} \rangle \langle e _ {\beta}, y ^ {\beta^ {\prime}} \rangle = c _ {\alpha \beta \gamma} \end{array}
\]

and hence

\[
e _ {\alpha} * e _ {\beta} = \sum_ {\gamma} c _ {\alpha \beta \gamma} e _ {\gamma}.\tag{11}
\]

(Formula (10) then expresses associativity on U(G).)

In particular, since L(G) is stable under the bracket.

\[
[ e _ {i}, e _ {j} ] = \sum_ {k} (c _ {i j k} - c _ {j i k}) e _ {k}.\tag{12}
\]

The constants of structure of \(\mathbf{L}(\mathbf{G})\) relative to the basis \((e_{1},\ldots,e_{n})\) are therefore the \(c_{ijk}-c_{jik}\) . In other words, canonically identifying \(\mathbf{L}(\mathbf{G})\) with \(K^{n}\) , we obtain:

\[
[ x, y ] = \mathrm{B} (x, y) - \mathrm{B} (y, x).\tag{13}
\]

PROPOSITION 1.

\begin{enumerate}
\def\labelenumi{(\roman{enumi})}
\tightlist
\item
\end{enumerate}

\[
x ^ {[ - 1 ]} \equiv - x + \mathrm{B} (x, x) \mod \deg 3\tag{i}
\]

\[
x. y. x ^ {[ - 1 ]} \equiv y + [ x, y ] \quad \text {   mod   } \deg 3\tag{ii}
\]

\[
y ^ {[ - 1 ]}. x. y \equiv x + [ x, y ] \quad \text {   mod   } \deg 3\tag{iii}
\]

\[
x ^ {[ - 1 ]}. y ^ {[ - 1 ]}. x. y \equiv [ x, y ] \quad \text {   mod   } \deg 3\tag{iv}
\]

\[
x. y. x ^ {[ - 1 ]}. y ^ {[ - 1 ]} \equiv [ x, y ] \quad \text {   mod   } \deg 3\tag{v}.
\]

(In (i), \(x^{[-1]}\) of course represents the expansion of the function \(x \mapsto x^{[-1]}\) as an integral series about the origin; the other formulae can be interpreted analogously.)

Let \(g_{1}\) and \(g_{2}\) be the homogeneous components of \(x^{[-1]}\) of degrees 1 and 2. Then

\[
\begin{array}{r l} 0 = x. x ^ {[ - 1 ]} \\ \equiv x + g _ {1} (x) + \mathrm{B} (x, g _ {1} (x)) & \text {   mod   } \deg 2 \quad \text {(by (6))} \\ \equiv x + g _ {1} (x) & \text {   mod   } \deg 2 \end{array}
\]

and hence \(g_{1}(x) = -x\) . Then

\[
\begin{array}{r l} 0 & = x. x ^ {[ - 1 ]} \\ & \equiv x + (- x + g _ {2} (x)) + B (x, - x + g _ {2} (x)) \mod \deg 3 \\ & \equiv g _ {2} (x) - B (x, x) \mod \deg 3 \end{array}
\]

and hence \(g_{2}(x) = \mathrm{B}(x, x)\) . This proves (i). Then, using (i),

\[
\begin{array}{l l} x. y. x ^ {[ - 1 ]} \equiv (x + y + \mathrm{B} (x, y)) \cdot (- x + \mathrm{B} (x, x)) & \text {mod} \deg 3 \\ \equiv x + y + \mathrm{B} (x, y) + (- x + \mathrm{B} (x, x)) + \mathrm{B} (x + y, - x) & \text {mod} \deg 3 \\ \equiv y + \mathrm{B} (x, y) - \mathrm{B} (y, x) & \text {mod} \deg 3 \\ \equiv y + [ x, y ] \qquad \qquad \qquad \qquad \qquad \qquad \qquad \qquad \qquad \qquad \qquad \qquad \qquad \qquad \qquad \qquad \qquad \qquad \qquad \qquad \qquad \qquad \qquad \qquad \qquad \qquad \qquad \qquad \qquad \qquad \qquad \qquad \qquad \qquad \qquad \qquad \qquad \\ & (\text {by (13)}) \end{array}
\]

whence (ii). The proof of (iii) is analogous. Combining (i) and (iii), we obtain

\[
\begin{array}{r l r} x ^ {[ - 1 ]}. y ^ {[ - 1 ]}. x. y & \equiv (- x + \mathrm{B} (x, x)). (x + [ x, y ]) & \bmod \deg 3 \\ & \equiv - x + \mathrm{B} (x, x) + x + [ x, y ] + \mathrm{B} (- x, x) & \bmod \deg 3 \\ & \equiv [ x, y ] & \bmod \deg 3 \end{array}
\]

whence (iv). The proof of (v) is analogous.

